\documentclass{article}
\usepackage{macros}
\usepackage{tkz-berge}
\usepackage[portuguese]{babel}
\usepackage[backend=biber, style=numeric, sorting=none]{biblatex}\addbibresource{referencias_grafos.bib}

\DeclareMathOperator{\tr}{tr}
\DeclareMathOperator{\spect}{spect}

\begin{document}

\begin{center}
    \large UNIVERSIDADE FEDERAL DO RIO GRANDE DO SUL \\
    \large INSTITUTO DE MATEMÁTICA E ESTATÍSTICA \\
    \vspace{4cm}
    \huge TÉCNICAS COMBINATÓRIAS E ANALÍTICAS EM TEORIA DE GRAFOS \\
    \vspace{1cm}
    \Large RELATÓRIO DE PROJETO DE INICIAÇÃO CIENTÍFICA \\
    \vspace{4cm}

    \large\begin{tabular}{rl}
        Acadêmico: & LEONARDO KAWAI SUAREZ \\
        Orientador: & Prof. Dr. VILMAR TREVISAN \\
    \end{tabular}

    \vfill
    \large Porto Alegre, RS \\
    \large \today
\end{center}

\newpage
\tableofcontents

\section{Preliminares sobre Teoria dos Grafos}

\begin{defn}
    Seja \(G=G(V,E)\) um grafo. Quando \(G'=G'(V',E')\) é um grafo satisfazendo \(V' \subseteq V\) e \(E' \subseteq E\), dizemos que \(G'\) é um \emph{subgrafo} de \(G\). Se \(G'\) é um subgrafo de \(G\) e \(E' = \{\{u,v\} \in E : u,v \in V'\}\), então \(G'\) é um \emph{subgrafo induzido} de \(G\).
\end{defn}

A seguir apresentamos alguns tipos especiais de grafos.

\begin{itemize}
    \item Um \emph{grafo regular de grau \(k\) ou \(k\)-regular} é um grafo \(G\) tal que \(\deg(v) = k\) para todo \(v \in V(G)\).
    \item Um \emph{grafo completo} é um grafo \(G\) tal que \(\{u,v\} \in E(G)\) para todo \(u,v \in V(G)\), \(u \neq v\). O grafo completo com \(n\) vértices é denotado por \(K_n\).
    \item Um \emph{grafo \(k\)-partido} é um grafo \(G=G(V,E)\) tal que \(V(G)\) pode ser particionado em \(k\) conjuntos independentes. O grafo \(k\)-partido completo é o grafo \(G\) tal que \(V(G)\) pode ser particionado em \(k\) conjuntos independentes e \(\{u,v\} \in E(G)\) para todo \(u,v \in V(G)\), \(u \neq v\), tais que \(u\) e \(v\) estão em conjuntos independentes distintos. O grafo bipartido completo com partição de tamanho \(m\) e \(n\) é denotado por \(K_{m,n}\). O grafo \(K_{1,n}\) é chamado de \emph{estrela}.
    \item Uma sequência finita de vértices \(v_1, v_2, \ldots, v_k\) de um grafo \(G(V,E)\) é uma \emph{cadeia} (\emph{walk}) de \(v_1\) a \(v_k\) se \(\{v_i, v_{i+1}\} \in E\) para todo \(i\), \(1 \leq i \leq k-1\). Uma \emph{cadeia fechada} (respectivamente, \emph{cadeia aberta}) é uma cadeia \(v_1, v_2, \ldots, v_k\) tal que \(v_1 = v_k\) (respectivamente, \(v_1 \neq v_k\)). Se todos os vértices de uma cadeia \(v_1, v_2, \ldots, v_k\) são distintos, então a cadeia é um \emph{caminho}. Um caminho fechado é denominado \emph{ciclo}. O \emph{comprimento} de um caminho ou de um ciclo é o número de arestas que ele contém. O caminho e o ciclo com \(n\) vértices são denotados por \(P_n\) e \(C_n\), respectivamente.
    \item Um grafo \(G\) é \emph{conexo} se para todo \(u,v \in V(G)\) existe uma cadeia de \(u\) a \(v\). Caso contrário, \(G\) é \emph{desconexo}. Se \(G\) é desconexo, dizemos que \(G' \subseteq G\) é uma \emph{componente conexa} de \(G\) se \(G'\) é um subgrafo induzido de \(G\) que é conexo e não existe um subgrafo induzido \(G''\) de \(G\) tal que \(G' \subset G'' \subseteq G\) e \(G''\) é conexo.
    \item Uma \emph{árvore} é um grafo conexo e acíclico. Um grafo desconexo acíclico é chamado de \emph{floresta}.
    \item Um \emph{grafo-linha} de um grafo \(G\) é o grafo \(L(G)\) cujo conjunto de vértices é \(E(G)\) e tais que \(\{e,f\} \in E(L(G))\) se e somente se \(e\) e \(f\) têm um vértice em comum em \(G\).
    \item O \emph{grafo complementar} de um grafo \(G\) é o grafo \(\overline{G}\) tal que \(V(\overline{G}) = V(G)\) e \(\{u,v\} \in E(\overline{G})\) se e somente se \(\{u,v\} \notin E(G)\).
    \item Quando \(V_1 \cap V_2 \neq \emptyset\), o \emph{grafo união} (ou \emph{soma}) de \(G_1\) e \(G_2\) é o grafo \(G_1 \cup G_2\), também denotado \(G_1 + G_2\), tal que \(V(G_1 \cup G_2) = V(G_1) \cup V(G_2)\) e \(E(G_1 \cup G_2) = E(G_1) \cup E(G_2)\). Para um grafo conexo \(G\), escrevemos \(kG\) para o grafo união de \(k\) cópias de \(G\).
    \item O \emph{produto cartesiano} de \(G_1\) e \(G_2\) é o grafo \(G_1 \square G_2\) tal que \(V(G_1 \square G_2) = V(G_1) \times V(G_2)\) e \(\{(u_1, u_2), (v_1, v_2)\} \in E(G_1 \square G_2)\) se e somente se \(\{u_1, v_1\} \in E(G_1)\) e \(u_2 = v_2\), ou \(u_1 = v_1\) e \(\{u_2, v_2\} \in E(G_2)\).
    \item O \emph{produto completo} (ou \emph{join}) de \(G_1\) e \(G_2\) é o grafo \(G_1 \vee G_2\) tal que \(V(G_1 \vee G_2) = V(G_1) \cup V(G_2)\) e \(E(G_1 \vee G_2) = E(G_1) \cup E(G_2) \cup \{\{u,v\} : u \in V(G_1), v \in V(G_2)\}\).
\end{itemize}

\begin{defn}
    Seja \(G=G(V,E)\) um grafo. O \emph{grau mínimo} de \(G\) é o número \(\delta(G) = \min\{\deg(v) : v \in V\}\), e o \emph{grau máximo} de \(G\) é o número \(\Delta(G) = \max\{\deg(v) : v \in V\}\). O \emph{grau médio} de \(G\) é o número \(\overline{\deg}(G) = \frac{1}{|V|} \sum_{v \in V} \deg(v)\).
\end{defn}

\begin{defn}
    Seja \(G = G(V,E)\) um grafo conexo. Se \(x, y \in V\), chamamos \emph{distância} de \(x\) a \(y\) e denotamos por \(d(x,y)\) ao mínimo dos comprimentos dos caminhos que ligam \(x\) a \(y\). O máximo das distâncias entre dois vértices quaisquer de \(G\) é chamado de \emph{diâmetro} de \(G\) e denotado por \(\operatorname{diam}(G)\). Quando \(G\) é desconexo, escrevemos \(\operatorname{diam}(G) = \infty\).
\end{defn}

\section{Matriz de adjacência}

\begin{defn}
    Seja \(G = G(V,E)\) um grafo com \(n\) vértices. A \emph{matriz de adjacência} de \(G\) é a matriz \(A(G) = (a_{ij})\) de ordem \(n\) definida por
    \[
        a_{ij} \defeq \begin{cases}
            1, & \text{se } \{v_i, v_j\} \in E\, \\
            0, & \text{caso contrário}.
        \end{cases}
    \]
\end{defn}

\subsection{Polinômio característico e espectro de um grafo}

Sendo \(A(G)\) uma matriz real simétrica, todos seus autovalores são reais.\footnote{Se \(T \in \mathcal{L}(V)\) é auto-adjunto (ou hermitiano) então todo autovalor de \(T\) é real.} Como \(\tr A(G) = 0\), a soma dos autovalores de \(A(G)\) é nula.

\begin{defn}
    O \emph{polinômio característico} de um grafo \(G\) é o polinômio característico da matriz de adjacência \(A(G)\), denotado por \(p_G(\lambda) = \det(\lambda I - A)\); \(\lambda\) é dito um \emph{autovalor do grafo} \(G\) quando \(p_G(\lambda) = 0\). Se \(A(G)\) possui \(n\) autovalores distintos \(\lambda_1, \ldots, \lambda_n\) com multiplicidades \(m_1, \ldots, m_n\), respectivamente, então o \emph{espectro} de \(G\), denotado por \(\spect(G)\), é a matriz \(2 \times n\) dada por
\[
    \spect(G) \defeq \begin{bmatrix}
        \lambda_1 & \lambda_2 & \cdots & \lambda_n \\
        m_1 & m_2 & \cdots & m_n
    \end{bmatrix}.
\]
O maior autovalor \(\lambda_1\) de \(G\) é chamado de \emph{índice} de \(G\). Para um inteiro \(k\geq0\), o \(k\)-ésimo \emph{momento espectral} de \(G\) é \(s_k = \sum_{i=1}^{k} \lambda_i^k\).
\end{defn}

Os autovalores de um grafo são independentes da ordenação de seus vértices, pois matrizes similares possuem o mesmo polinômio característico: se os vértices são permutados, obtemos uma matriz de adjacência \(A'\) que é similar a \(A\) por meio de uma matriz de permutação \(P\), ou seja, \(A' = P^T A P\).

\begin{prop}
    Seja \(G\) um grafo com \(n\) vértices e \(m\) arestas e seja
    \begin{equation}
        p_G(\lambda) = \lambda^n + a_1 \lambda^{n-1} + a_2 \lambda^{n-2} + \cdots + a_{n-1} \lambda + a_n
    \end{equation}
    o polinômio característico de \(G\). Então \(a_1 = 0\), \(a_2 = -m\) e \(a_n = -2t\), onde \(t\) é o número de triângulos em \(G\).
\end{prop}

\begin{proof}
    Da teoria de matrizes, temos que para todo \(k\), \(1 \leq k \leq n\), vale 
    \[
        a_k = \sum (\text{todos os \(k \times k\) menores principais}\footnotemark).
    \]\footnotetext{Um \emph{menor principal} de \(A(G)\) com \(k\) linhas e \(k\) colunas é o determinante de qualquer \emph{submatriz principal} com \(k\) linhas e \(k\) colunas de \(A(G)\), isto é, de qualquer matriz obtida de \(A(G)\) pela retirada de um subconjunto de \(n-k\) linhas e do correspondente subconjunto de colunas.}
\end{proof}

Os autovalores de \(A\) são os números reais que satisfazem \(Ax = \lambda x\) para algum vetor \(x \neq 0\). Tal vetor \(x\) é chamado de \emph{autovetor} de \(A\) associado ao autovalor \(\lambda\). A relação \(Ax = \lambda x\) pode ser entendida da seguinte forma: se \(x = (x_1, x_2, \ldots, x_n)^T\) então

\begin{equation}
    \lambda x_u = \sum_{v \sim u} x_v, \quad u = 1, 2, \ldots, n,
\end{equation} onde a soma é feita sobre os vizinhos \(v\) de \(u\). Assim, o autovetor \(x\) é uma função que atribui um número real a cada vértice de \(G\) de tal forma que para cada vértice \(u\) o valor \(\lambda x_u\) é igual à soma dos valores atribuídos aos vizinhos de \(u\).

\begin{prop}
    Se um grafo \(G\) tem grau máximo \(\Delta(G)\) então todo autovalor \(\lambda\) de \(G\) satisfaz \(|\lambda| \leq \Delta(G)\).
\end{prop}

\begin{proof}
    
\end{proof}

Se \(\lambda\) é um autovalor de \(A\) então o conjunto \(\{x \in \mathbb{R}^n : Ax = \lambda x\}\) é um subespaço vetorial de \(\mathbb{R}^n\) chamado de \emph{autoespaço} de \(\lambda\) e denotado por \(\mathcal{E}(\lambda)\) ou \(\mathcal{E}_A(\lambda)\). Tais autoespaços são chamados de \emph{autoespaços de \(G\)}. Como \(A\) é uma matriz real simétrica, o teorema espectral assegura que \(A\) pode ser diagonalizada por uma matriz ortogonal, ou seja, existe uma base ortonormal de \(\mathbb{R}^n\) formada por autovetores de \(A\). Então os autoespaços de \(G\) são dois a dois ortogonais e a soma direta de todos os autoespaços de \(G\) é igual a \(\mathbb{R}^n\). Além disso, a dimensão de \(\mathcal{E}(\lambda)\) é igual à multiplicidade de \(\lambda\) como raiz de \(p_G(\lambda)\), isto é, a multiplicidade geométrica de \(\lambda\) é igual à sua multiplicidade algébrica. Por isso, nos referimos apenas à \emph{multiplicidade} de \(\lambda\). Um autovalor \emph{simples} é aquele cuja multiplicidade é \(1\).

\begin{prop}
    O número de cadeias de comprimento \(l\) ligando os vértices \(v_i\) e \(v_j\) em um grafo \(G\) é dado por \((A^l)_{ij}\).
\end{prop}

\begin{proof}
    Provamos por indução em \(l\). Trivial para \(l=1\). Suponha que vale para \(l=n\) e nota que a quantidade de cadeias de comprimento \(l+1\) ligando \(v_i\) a \(v_j\) é igual à quantidade de cadeias de comprimento \(l\) ligando \(v_i\) a um vizinho \(v_h\) de \(v_j\). Então
        \[
            \sum_{\{v_i\sim v_h\} \in E} (A^l)_{ih} = \sum_{h=1}^n (A^l)_{ih}A_{hj} = (A^{l+1})_{ij}.
        \]
    \end{proof}
\end{proof}

Podemos usar o resultado acima para relacionar o número de cadeias fechadas de um grafo e as somas de potências de seus autovalores.

\begin{cor}
    Seja \(G\) um grafo com \(n\) vértices e \(m\) arestas. Se \(T_l\) é o número total de cadeias fechadas de comprimento \(l\) em \(G\), então \(T_l = \tr A^l = \sum_{i=1}^n \lambda_i^l\), onde \(\lambda_1, \ldots, \lambda_n\) são os autovalores de \(G\). Em particular,
    \begin{enumerate}
        \item \(T_2 = 2m\), isto é, o número de cadeias fechadas de comprimento 2 em \(G\) é igual a duas vezes o número de arestas de \(G\);
        \item se \(G\) é um grafo regular de grau \(k\) então \(T_2 = kn\), pois \(kn=2m\);
        \item a soma dos cubos dos autovalores de \(G\) é igual a seis vezes o número de triângulos em \(G\), ou seja, \(T_3 = \sum_{i=1}^n \lambda_i^3 = 6t\).
    \end{enumerate}
\end{cor}

\subsection{Alguns resultados e aplicações}

\begin{thm} [teorema do entrelaçamento de Cauchy para a matriz de adjacência]
    Seja \(G = G(V,E)\) um grafo com \(n\) vértices e \(H = H(V',E')\) um subgrafo induzido de \(G\) com \(m\) vértices. Se \(\lambda_1 \geq \cdots \geq \lambda_n\) são os autovalores de \(G\) e \(\mu_1 \geq \cdots \geq \mu_m\) são os autovalores de \(H\), então para todo \(i\), \(1 \leq i \leq m\), vale
    \[
        \lambda_i \geq \mu_i \geq \lambda_{n-m+i}.
    \]
    Em particular, se \(v \in V\) e \(G-v\) é o subgrafo induzido de \(G\) obtido pela remoção do vértice \(v\), então para todo \(i\), \(1 \leq i \leq n-1\), vale
    \[
        \lambda_i \geq \mu_i \geq \lambda_{n-1+i}.
    \]
\end{thm}

\begin{defn}
    O \emph{raio espectral} de um grafo \(G\), denotado por \(\rho(G)\), é o número real não negativo \(\rho(G) \defeq \max\{|\lambda| : \lambda \in \spect(G)\}\).
\end{defn}

\begin{defn}
    A \emph{centralidade de autovetor} do vértice \(v_i\) em um grafo \(G\) é a \(i\)-ésima coordenada \(x_i\) do autovetor unitário positivo \(x = [x_1 \; \ldots \; x_n]^T\) associado ao autovalor principal \(\lambda_1\) de \(G\), ou seja, é o número
    \begin{equation}
        x_i = \frac{1}{\lambda_1} \sum_{j=1}^n a_{ij} x_j,
    \end{equation}
    onde \(A(G) = (a_{ij})\) é a matriz de adjacência de \(G\).
\end{defn}

\begin{prop}
    Se \(G\) é um grafo conexo então o número de autovalores distintos de \(G\) é \(> \operatorname{diam}(G)\).
\end{prop}

\begin{proof}
    Provamos por contradição. Sejam \(\lambda_1, \ldots, \lambda_t\) os autovalores distintos de \(G\). Como \(A\) é uma matriz real e simétrica, seu polinômio minimal possui \(t\) fatores lineares distintos e então vale
    \[
        (A-\lambda_1) \cdots (A-\lambda_t) = 0,
    \]
    o que mostra que \(A^t\) é combinação linear de \(A, A^2, \ldots A^{t-1}\). Suponha que \(t = \operatorname{diam}(G)\) e tome vértices \(u, v\) tal que \(d(u,v) = t\). Então \(A_{uv}^k = 0\) para todo \(k < t\) mas \(A_{uv}^t > 0\), o que é uma contradição pois \(A^t\) é combinação linear dos \(A^k\). Logo, \(t > \operatorname{diam}(G)\).
\end{proof}

\begin{prop}
    Se \(G\) é um grafo com maior autovalor \(\lambda_1\) então
    \begin{equation}
        \overline{\deg}(G) \leq \lambda_1 \leq \Delta(G).
    \end{equation}
    onde \(\overline{\deg}(G)\) é o grau médio de \(G\) e \(\Delta(G)\) é o grau máximo de \(G\).
\end{prop}

\begin{proof}
    
\end{proof}

\subsection{O espectro de certos tipos de grafos}

Encontrar o espectro de grafos com mais do que quatro vértices através do cálculo das raízes de seu polinômio característico é uma tarefa difícil que requer uso de programas de computação algébrica. Existem resultados que determinam o espectro de grafos obtidos a partir de outros por operações combinatórias, como o complementar, a união etc. Outros teoremas fornecem informações sobre apenas alguns autovalores. No geral, são conhecidos os espectros de poucas classes de grafos.

\begin{prop} \label{prop:polinomio_uniao}
    Se \(G\) é a união de dois grafos \(G_1\) e \(G_2\) então \(p_G(\lambda) = p_{G_1}(\lambda) p_{G_2}(\lambda)\). Em particular, se \(G_1, G_2, \ldots, G_r\) são as componentes conexas de \(G\) então \(p_G(\lambda) = \prod_{i=1}^r p_{G_i}(\lambda)\).
\end{prop}

\begin{proof}
    
\end{proof}

\begin{prop}
    Seja \(G\) um grafo regular de grau \(k\). Então
    \begin{enumerate}
        \item [a.] \(k\) é um autovalor de \(G\);
        \item [b.] \(G\) é conexo se e somente se \(k\) é um autovalor simples de \(G\);
        \item [c.] qualquer autovalor \(\lambda\) de \(G\) satisfaz \(|\lambda| \leq k\).
    \end{enumerate}
\end{prop}

\begin{proof}
    \begin{enumerate}
        \item [a.] Perceba que a soma das entradas da linha \(j\) da matriz de adjacência fornece o grau do vértice \(v_j\). Como \(G\) é \(k\)-regular, todas essas somas resultam em \(k\). O candidato natural a autovetor associado a \(k\) é o autovetor constante. De fato, tomando \(x=[1 \ldots 1]^T\) tem-se \(A1 = k1\), donde \(k\) é autovalor de \(G\).
        \item [b.] \((\Rightarrow)\) Seja \(x = [x_1 \; \ldots \; x_n]^T\) autovetor de \(G\) associado ao autovalor \(k\). Provamos que \(x\) é múltiplo escalar do vetor \(1\). Toma \(x_j = \max\{\vert x_i \vert; 1 \leq i \leq n\}\) e nota que \((Ax)_j = kx_j = \sum_{v_i \sim v_j} x_i\) (soma dos valores dos \(k\) vizinhos de \(x_j\)). Então para todo \(v_l\) adjacente a \(v_j\) vale
        \begin{equation*}
            k \vert x_j \vert = \vert \sum_{v_i \sim v_j} x_i \vert \leq \sum_{v_i \sim v_j} \vert x_i \vert \leq \vert x_l \vert + (k-1) \vert x_j \vert.
        \end{equation*} Subtraindo \((k-1) \vert x_k \vert\) de ambos os lados, vem \(\vert x_j \vert \leq \vert x_l \vert\). Dado que tomamos \(x_j\) como o valor absoluto máximo de \(x\), devemos ter \(x_l = x_j\). Ainda, a conexidade do grafo nos permite aplicar essa lógica a todos os seus vértices. Portanto \(x = x_j1\) e o autoespaço associado a \(k\) tem dimensão \(1\).
        
        \((\Leftarrow)\) Por outro lado, suponha que \(G\) é desconexo, com componentes conexas \(G_1, \ldots, G_m\). Como cada componente é um grafo conexo \(k\)-regular, \(k\) é autovalor simples para cada \(G_i\). Da proposição anterior, temos que \(p_G(\lambda) = p_{G_1}(\lambda)\cdots p_{G_m}(\lambda)\), logo \(k\) é autovalor de \(G\) de multiplicidade \(>1\).
        \item [c.] Seja \(y = [y_1 \; \ldots \; y_n]^T\) autovetor de \(G\) associado ao autovalor \(\lambda\). Defina \(y_j = \max\{\vert y_i \vert; 1 \leq i \leq n\}\) e escreva \((Ay)_j = \lambda y_j = \sum_{v_i \sim v_j} y_i\). Daí, tem-se 
        \begin{equation*}
            \vert \lambda \vert \vert y_j \vert = \left\vert \sum_{v_i \sim v_j} y_i \right\vert \leq \sum_{v_i \sim v_j} \vert y_i \vert \leq k \vert y_j \vert,
        \end{equation*} donde \(\lambda \leq k\).
    \end{enumerate}
\end{proof}

Geralmente, não se conhece todo o espectro de um grafo \(k\)-regular. Mas se conhecemos seus autovalores, podemos determinar os autovalores de seu grafo complementar.

\begin{prop}
    Seja \(G\) um grafo \(k\)-regular com \(n\) vértices e autovalores \(\lambda_1 = k, \lambda_2, \ldots, \lambda_n\). Então os autovalores de \(\overline{G}\) são \(n-1-k, -1-\lambda_2, \ldots, -1-\lambda_n\), respectivamente, associados aos mesmos autovetores de \(G\).
\end{prop}

\begin{proof}
    
\end{proof}

\begin{exmp}
    Se \(G\) é um grafo sem arestas com \(n\) vértices, então \(G\) é \(0\)-regular com \(p_G(\lambda) = \lambda^n\) e \[\spect(G) = \begin{bmatrix}
        0 \\
        n
    \end{bmatrix}.\]
\end{exmp}

\begin{exmp}
    O grafo completo com \(n\) vértices \(K_n\) é o complementar do grafo sem arestas. Temos então \(p_{K_n}(\lambda) = (\lambda - (n-1))(\lambda + 1)^{n-1}\) e \[\spect(K_n) = \begin{bmatrix}
        n-1 & -1 \\
        1 & n-1
    \end{bmatrix}.\]
\end{exmp}

\begin{exmp}
    Se \(G\) é o grafo formado pela união de \(s\) cópias de \(K_2\) então \(G\) é \(1\)-regular e tem \(2s\) vértices, com \(p_{K_2} (\lambda) = (\lambda - 1)(\lambda + 1)\) e por \ref{prop:polinomio_uniao}, vale \(p_G (\lambda) = (\lambda - 1)^s(\lambda + 1)^s\) e \[ \spect(G) = \begin{bmatrix}
        1 & -1 \\
        s & s
    \end{bmatrix}.\] O grafo complementar de \(G\) é chamado de \emph{hiperoctaedro} (ou \emph{cocktail party graph}) e denotado por \(H_s\). O espectro de \(H_s\) é \[\spect(H_s) = \begin{bmatrix}
        2s-1 & -1 \\
        1 & 2s-1
    \end{bmatrix}.\]
\end{exmp}

\begin{exmp}
    Seja \(G\) um grafo bipartido cujo conjunto de vértices é particionado em \(V_1\) e \(V_2\) com \(|V_1| = k_1\) e \(|V_2| = k_2\). Se rotulamos primeiro os vértices de \(V_1\) e depois os de \(V_2\), a matriz de adjacência de \(G\) tem a forma
\[
    A(G) = \begin{bmatrix}
        0 & B \\
        B^T & 0
    \end{bmatrix},
\] onde \(B\) é uma matriz de adjacência de dimensões \(k_1 \times k_2\). Afirmamos que \(\lambda\) é um autovalor de \(G\) se e somente se \(-\lambda\) é um autovalor de \(G\). De fato, se \(x = [x_1 \; x_2]^T\) é um autovetor de \(G\) associado a \(\lambda\), onde \(x_1\) e \(x_2\) são vetores de dimensões \(k_1\) e \(k_2\), respectivamente, então
\[
    \lambda x_1 = Bx_2 \quad \text{e} \quad \lambda x_2 = B^T x_1.
\] Daí, \(-\lambda x_1 = -Bx_2\) e \(-\lambda x_2 = -B^T x_1\), o que mostra que \(-\lambda\) é um autovalor de \(G\) associado ao autovetor \([-x_1 \; x_2]^T\). O espectro de \(G\) é então da forma
\[
    \spect(G) = \begin{bmatrix}
        \lambda_1 & \lambda_2 & \cdots & \lambda_r & -\lambda_1 & -\lambda_2 & \cdots & -\lambda_r \\
        m_1 & m_2 & \cdots & m_r & m_1 & m_2 & \cdots & m_r
    \end{bmatrix},
\] onde \(\lambda_1, \ldots, \lambda_r\) são os autovalores positivos de \(G\) e \(m_1, \ldots, m_r\) são suas respectivas multiplicidades.
\end{exmp}

\subsection{Isomorfismo de grafos}

\begin{defn}
    Dois grafos \(G_1\) e \(G_2\) são \emph{isomorfos} se existe uma bijeção \(f: V(G_1) \to V(G_2)\) tal que \(\{u,v\} \in E(G_1)\) se e somente se \(\{f(u), f(v)\} \in E(G_2)\). Neste caso, dizemos que \(f\) é um \emph{isomorfismo} entre \(G_1\) e \(G_2\).
\end{defn}

\begin{defn}
    Dois grafos \(G_1\) e \(G_2\) são \emph{cospectrais} quando têm os mesmos autovalores com as mesmas multiplicidades, ou seja, \(\spect(G_1) = \spect(G_2)\).
\end{defn}

Se dois grafos são isomorfos, eles possuem o mesmo espectro. Mas a recíproca não é verdadeira, em geral.

\begin{defn}
    Dizemos que um grafo \(G\) é \emph{determinado pelo espectro} se para todo grafo \(H\) cospectral a \(G\), \(H\) é isomorfo a \(G\).
\end{defn}

\begin{center}
    \begin{tikzpicture}
        \GraphInit[vstyle=Shade]
        \SetVertexMath
        \grCycle[prefix=v,RA=3]{4}
        \grCycle[prefix=w,RA=1.3]{4}
        \EdgeIdentity{v}{w}{4}
    \end{tikzpicture}
\end{center}

Algumas propriedades de grafos, como a conexidade e a existência de \(C_4\), não são determinadas por seus espectros. O resultado a seguir mostra que grafos sem arestas são caracterizados pelo seu espectro.

\begin{prop}
    Um grafo \(G\) possui um único autovalor se e somente se \(G\) é totalmente desconexo, ou seja, é um grafo sem arestas.
\end{prop}

\begin{proof}
    
\end{proof}

\begin{prop}
    Se \(G\) tem exatamente dois autovalores distintos \(\lambda_1 > \lambda_2\), com multiplicidades \(m_1\) e \(m_2\), respectivamente, então \(G\) é o grafo \(\lambda_1\)-regular formado por \(m_1\) cópias de \(K_{\lambda_1 + 1}\). Além disso, temos \(\lambda_2 = -1\) e \(m_2 = m_1\lambda_1\).
\end{prop}

\begin{proof}
    
\end{proof}

\subsection{Energia de grafos}

\begin{defn}
    A \emph{energia de um grafo} \(G\) com \(n\) vértices é definida por
\[
    \mathcal{E}(G) \defeq \sum_{i=1}^n |\lambda_i|,
\]
onde \(\lambda_1, \ldots, \lambda_n\) são os autovalores de \(G\).
\end{defn}

\section{Matriz de incidência}

\subsection{Matriz de incidência e grafo-linha de um grafo}

\section{Matriz Laplaciana}

\begin{defn}
    Seja \(D\) a \emph{matriz diagonal dos graus} dos vértices de um grafo \(G\) (isto é, \(D_{ii} = \deg(v_i)\) para todo \(i\)), e \(A(G)\) a matriz de adjacência de \(G\). A \emph{matriz Laplaciana} (ou \emph{Laplaciano}) de \(G\) é a matriz \(L(G) = D - A(G)\).
\end{defn}

\begin{defn}
    O \emph{espectro do Laplaciano} \(L(G)\) de um grafo \(G\), denotado por \(\xi(G)\), é a matriz-linha dos autovalores de \(L(G)\) ordenados de forma não crescente. Assim, se \(\mu_1 \geq \cdots \geq \mu_n\) são os autovalores de \(L(G)\), então \(\xi(G) = [\mu_1, \ldots, \mu_n]\).
\end{defn}

\begin{defn}
    Seja \(G\) um grafo. A \emph{matriz de incidência com respeito a uma orientação} de \(G\) é a matriz \(\beta(G)\) definida por
    \[
        \beta_{ij} \defeq \begin{cases}
            1, & \text{se } e_j = (v_i, w) \text{ para algum } w \in V(G), \\
            -1, & \text{se } e_j = (w, v_i) \text{ para algum } w \in V(G), \\
            0, & \text{caso contrário}.
        \end{cases}
    \]
\end{defn}

\subsection{Teorema da Matriz-Árvore}

\begin{defn}
    Uma \emph{árvore geradora} de um grafo \(G\) é um subgrafo de \(G\) que é uma árvore e contém todos os vértices de \(G\).
\end{defn}



\subsection{Conectividade algébrica}

\begin{defn}
    A \emph{conectividade algébrica} de um grafo \(G\) é o segundo menor autovalor do Laplaciano de \(G\), denotado por \(a(G) \defeq \mu_{n-1}(G)\). O \emph{índice do Laplaciano} de \(G\) é o maior autovalor do Laplaciano de \(G\), denotado por \(\mu_1(G)\).
\end{defn}

\begin{prop}

\end{prop}

\begin{cor}

\end{cor}

\begin{defn}
    A \emph{conectividade de vértices} de um grafo \(G\), denotada por \(\kappa(G)\), é o número mínimo de vértices que devem ser removidos para desconectar \(G\) ou torná-lo trivial. A \emph{conectividade de arestas} de \(G\), denotada por \(\kappa'(G)\), é o número mínimo de arestas que devem ser removidas para desconectar \(G\) ou torná-lo trivial.
\end{defn}

\subsection{Autovalores e o problema do corte maximal}



\subsection{Aplicações}

\section{Matriz Laplaciana sem sinal}

\begin{defn}
    A \emph{matriz Laplaciana sem sinal} de um grafo \(G\) é a matriz \(Q(G) = D + A(G)\), onde \(D\) é a matriz diagonal dos graus dos vértices de \(G\) e \(A(G)\) é a matriz de adjacência de \(G\).
    O \emph{polinômio característico da matriz Laplaciana sem sinal} é denotado por \(q_G(\lambda) = \det(\lambda I - Q(G))\), e os autovalores de \(Q(G)\) são as raízes de \(q_G(\lambda)\), denotados por \(q_1 \geq \cdots \geq q_n\), sendo \(q_1\) o \emph{índice de \(Q(G)\)}.
\end{defn}

\subsection{\(Q\)-espectro da união, produto cartesiano e join}
\subsection{Grafos \(Q\)-integrais}
\subsection{Coespectralidade e aplicações}

\section{Algoritmos de Localização de Autovalores}

\begin{defn}
    Diz-se que duas matrizes \(R\) e \(S\) são \emph{congruentes} quando existe uma matriz invertível \(P\) tal que \(R = P^T S P\).
\end{defn}

\begin{thm} [lei de inércia de Sylvester]
    Duas matrizes simétricas reais são congruentes se e somente se possuem o mesmo número de autovalores negativos e o mesmo número de autovalores positivos.
\end{thm}

\subsection{Adjacência de árvores}
\subsection{Aplicações para árvores}
\subsection{Adjacências de grafos threshold}

\begin{appendix}

\section{Resultados de Teoria Matricial}

\subsection{Princípio de Rayleigh}

Para os resultados a seguir, sejam \(A_{n \times n}\) uma matriz real simétrica com autovalores \(\lambda_1 \geq \ldots \geq \lambda_n\), \(\langle \cdot, \cdot \rangle\) o produto interno canônico em \(\R^n\) e \(\|\cdot\|\) a norma euclidiana induzida por esse produto interno.



\begin{prop}
    \begin{equation}
        \lambda_1 = \max_{\Vert x \Vert = 1} \langle Ax, x \rangle.
    \end{equation}
\end{prop}

\begin{cor}
    \begin{equation}
        \lambda_1 = \min_{x \neq 0} \frac{\langle Ax, x \rangle}{\langle x, x \rangle}.
    \end{equation}
    e o valor máximo é atingido por um autovetor unitário associado a \(\lambda_1\).
\end{cor}

\begin{prop}
    Para todo \(k\), \(1 \leq k \leq n\), vale
    \begin{equation}
        \lambda_k = \max{\langle Ax, x \rangle : \Vert x \Vert = 1, \langle x, v_i \rangle = 0 \quad\forall i, 1 \leq i < k},
    \end{equation}
    sendo o valor máximo atingido por um autovetor unitário associado a \(\lambda_k\).
\end{prop}

\begin{cor}
    Para todo \(k\), \(1 \leq k \leq n\), vale
    \begin{equation}
        \lambda_k = \max\left\{\frac{\langle Ax, x \rangle}{\langle x, x \rangle} : x \neq 0, \langle x, v_i \rangle = 0 \quad\forall i, 1 \leq i < k\right\},
    \end{equation}
    e o valor máximo é atingido por um autovetor associado a \(\lambda_k\).
\end{cor}

\begin{cor}
    Para todo \(k\), \(1 \leq k \leq n\), vale
    \begin{equation}
        \lambda_k = \min\left\{\frac{\langle Ax, x \rangle}{\langle x, x \rangle} : x \neq 0, \langle x, v_i \rangle = 0 \quad\forall i, k < i \leq n\right\},
    \end{equation}
    e o valor mínimo é atingido por um autovetor associado a \(\lambda_k\).
\end{cor}

\subsection{Teorema de Entrelaçamento de Cauchy}

\begin{thm}
    Seja \(M_{n \times n}\) uma matriz real tal que \(M^TM = I\). Se os autovalores de \(A\) são \(\lambda_1 \geq \cdots \geq \lambda_n\) e os autovalores de \(M^T A M\) são \(\mu_1 \geq \cdots \geq \mu_m\), então para todo \(i\), \(1 \leq i \leq m\), vale
    \begin{equation}
        \lambda_i \geq \mu_i \geq \lambda_{n-m+i}.
    \end{equation}
\end{thm}

\subsection{Teorema de Perron-Frobenius}

O Teorema de Perron-Frobenius é um resultado fundamental da teoria de matrizes que tem diversas aplicações em sistemas dinâmicos, probabilidade e teoria de grafos. Para enunciar o teorema, precisamos de algumas definições.

\begin{defn}
    \(A \in \R^{m \times n}\) é chamada de \emph{matriz não negativa} quando cada \(a_{ij} \geq 0\), e escreve-se \(A \geq 0\). Em geral, \(A \geq B\) significa que \(a_{ij} \geq b_{ij}\) para todo \(i\) e \(j\). Similarmente, \(A\) é uma \emph{matriz positiva} quando cada \(a_{ij} > 0\), e escreve-se \(A > 0\). Mais geralmente, \(A > B\) significa que \(a_{ij} > b_{ij}\) para todo \(i\) e \(j\).    
\end{defn}

\begin{defn}
    Uma matriz simétrica \(A \in \R^{n \times n}\) é \emph{irredutível} se não existe uma matriz de permutação \(P\) tal que
    \[
        P^TAP = \begin{bmatrix}
            X & Z \\
            0 & Y
        \end{bmatrix},
    \]
    onde \(X\) e \(Y\) são matrizes quadradas. Caso contrário, \(A\) é \emph{redutível}. \(P^TAP\) é dita uma \emph{permutação simétrica} de \(A\), cujo efeito é permutar as linhas e as colunas de \(A\) da mesma forma.
\end{defn}

\begin{defn}
    O \emph{grafo} \(\mathcal{G}(A)\) de \(A\) é definido como o grafo direcionado em \(n\) vértices \(v_1, \ldots, v_n\) tal que \(\{v_i, v_j\}\) é uma aresta de \(\mathcal{G}(A)\) se e somente se \(a_{ij} \neq 0\). \(\mathcal{G}(A)\) é \emph{fortemente conexo} se para todo \(i\) e \(j\), \(1 \leq i,j \leq n\), existe um caminho de \(v_i\) a \(v_j\) em \(\mathcal{G}(A)\).
\end{defn}

\begin{prop}
    \(A\) é irredutível se e somente se \(\mathcal{G}(A)\) é fortemente conexo.
\end{prop}

\begin{lem}
    Se \(A \in \R^{n \times n}\) é não negativa e irredutível então \((I+A)^{n-1}>0\), ou equivalentemente, \(I + A + A^2 + \cdots + A^{n-1} > 0\).
\end{lem}

\begin{proof}
    
\end{proof}

\begin{prop}
    Se \(G\) é um grafo conexo com \(n\) vértices então \(A(G)\) é irredutível.
\end{prop}

\begin{proof}
    
\end{proof}

\begin{thm} [Teorema de Perron-Frobenius]
    Seja \(A_{n \times n}\) uma matriz não negativa e irredutível e com autovalores \(\lambda_1 \geq \ldots \geq \lambda_n\). Então
    \begin{enumerate}
        \item [a.] \(\lambda_1 > 0\) é um autovalor simples de \(A\) com um autovetor associado positivo;
        \item [b.] \(\lambda_1 > \lambda_2\);
        \item [c.] \(\vert \lambda_i \vert \leq \lambda_1\) para todo \(i\), \(1 \leq i \leq n\).
    \end{enumerate}
\end{thm}

\end{appendix}


\newpage

\printbibliography[title={Referências}]

\end{document}